\scp{Organisation of the book}{01-04}
\Crf{ch:IntCom} identifies the purpose
and audience of this study,
along with an overview of previous studies that have contributed to our collective understanding of the languages within the target region as a whole.

\Crf{ch:LayAdd} provides background to the complexities of the region,
the geography, prehistory, early trade
and migrations, and methodological issues of unravelling the transcription of old data sources
and similar issues that are a necessary backdrop to understanding the linguistic context.

Chapters \ref{ch:BL}--\ref{ch:Aru} examine the data
and identify the innovations that define four subgroups along the southern islands of our target region,
moving from west to \ili{east}.

\Crf{ch:CM} introduces the complexities of issues in geographical central Maluku (the northern chain of islands within our target region).
This interlude is necessary because much has been assumed or claimed,
yet there is very little agreement among scholars about how these languages of central Maluku subgroup,
and how they are or are not related to each other.

Chapters \ref{ch:Tal}--\ref{ch:Ban} examine the data
and identify the innovations that define five subgroups in geographical central Maluku,
stretching from northwest to south\ili{east}.

\Crf{ch:OthSub} examines the pros
and cons of various arguments that might be useful for linking various of our subgroups together at a higher level.
Several of these have been proposed in the literature for subgrouping within the region.

\Crf{ch:Mar} reexamines,
expands the data,
and refines arguments
and glosses related to marsupials,
and their implications for subgrouping.
\citet[727]{Blust2013} observed,
“the distribution of cognate sets for marsupial mammals must be considered the
jewel in the crown for this [\ili{Central} Eastern Malayo-Polynesian] hypothesis”.

\Crf{ch:MorLes} looks at the distribution of various lexical items
and other features that have been noted to be of significance for the target region,
along with higher level implications of our conclusions,
both for our understanding of the prehistory
and classification of the region,
and also for methodology.
Mechanisms are proposed for ways in which many lexical items
and other features can have the uneven distributions that we find within the region.
Several possible incursion
and migration scenarios are discussed.
A glimpse of what PMP might look like based on our findings is also given.

Appendix \ref{ch:LanLis} lists our data sources for each language.
Appendix \ref{ch:BanTer} lists words for ‘banana’ for 389 Austronesian
and Papuan languages,
along with their sources.
Appendix \ref{ch:SurMar} lists marsupial terms for 270 Austronesian
and Papuan languages,
along with their sources.
Appendix \ref{ch:ComWor} provides a comparative word list designed to help investigators
collect data which will be useful for comparative
purposes for languages of eastern Indonesia
and Timor-Leste.
All these appendices can be downloaded in .csv format
from XXX.

The chapters are not all of the same nature.
Chapters \ref{ch:BL}--\ref{ch:Aru},
and \ref{ch:SB}--\ref{ch:Ban} focus on specific details of reconstruction,
sound changes, and the internal structure of different sub-nodes.
We expect that these chapters are likely to be of more interest to historical linguists than to other specialists.
Chapters \ref{ch:IntCom}, \ref{ch:LayAdd}, \ref{ch:CM},
\ref{ch:OthSub}--\ref{ch:MorLes} are slightly more “general”
and deal with broader
and higher level issues.
Chapters \ref{ch:LayAdd} and \ref{ch:CM} provide important background
and context for understanding the region,
while chapters \ref{ch:OthSub}--\ref{ch:MorLes} summarise
and draw the discussion together.
It is our expectation that these broader chapters will be of more interest to those who are not historical linguists
and/or those unfamiliar with the Austronesian world.